\subsection{Four arc gains}

Theorems~\ref{thm:posprod} and~\ref{thm:posprodq} put one gain on one arc. This subsection allows every arc of the two-rail digraph its own gain and asks the same question. The answer depends on the four gains only through one quotient.

\paragraph{Setting.} $D$ is the two-rail digraph of the line or of the ring $\mathbb{Z}_N$ (vertices $(x,R)$, $(x,L)$; a right step from either to $(x+1,R)$, a left step from either to $(x-1,L)$), with gains in a finite abelian group $A$: $a$ on the arcs from $(x,R)$ to $(x+1,R)$, $b$ from $(x,R)$ to $(x-1,L)$, $c$ from $(x,L)$ to $(x+1,R)$, $d$ from $(x,L)$ to $(x-1,L)$. Put
\[
\delta=ad\,(bc)^{-1},\qquad r=\operatorname{ord}\delta,\qquad H=\langle a,d,bc\rangle=\langle a,bc,\delta\rangle,\qquad M=N/\gcd(N,2),\qquad \Delta_N=\langle\delta,(a/d)^{M}\rangle .
\]
For $X\subseteq A$ let $\mathcal{M}(X)$ be the multisets of elements of $X$, written additively, and $T_g$ the translation of a multiset by $g$. Events are sets of walks from $(0,R)$, position product events are as in the subsection on them above, and $m_S(y)\in\mathcal{M}(A)$ is the multiset of gains of the walks of $S$ that end at $y$. The rule adds $F_{(x,R)}(m_S(x,R))+F_{(x,L)}(m_S(x,L))$ over the sites, with $F(\emptyset)=0$. For one arc of gain $h$ on the arcs from $(x,L)$ to $(x-1,L)$, that is $(a,b,c,d)=(1,1,1,h)$, this is the setting of Theorem~\ref{thm:posprodq}; we write $D_h$ for that carrier.

\begin{lemma}[gains at one length]\label{lem:fgcoset}
Let $w$ be a walk from $(0,R)$ with $\rho$ right and $\lambda$ left steps, and $n$ arcs from rail $L$ to rail $R$. Its gain is $a^{\rho}d^{\lambda}\delta^{-n}$ if it ends on $R$ and $a^{\rho}d^{\lambda-1}b\,\delta^{-n}$ if it ends on $L$. Consequently, on the line, the gains of the walks of length $T$ that end at $(x,R)$ lie in one coset $\kappa\langle\delta\rangle$ with $\kappa=a^{x}(bc)^{(T-x)/2}\in H$, and those that end at $(x,L)$ lie in $\kappa\,b\,d^{-1}\langle\delta\rangle$; on $\mathbb{Z}_N$ the same holds with $\Delta_N$ in place of $\langle\delta\rangle$.
\end{lemma}
\begin{proof}
The rail after a step is $R$ exactly when the step is to the right. With $n_{XY}$ the number of arcs from rail $X$ to rail $Y$, $n_{RR}+n_{LR}=\rho$ and $n_{RL}+n_{LL}=\lambda$; the changes of rail alternate and the first is from $R$, so $n_{RL}=n_{LR}=n$ for a walk that ends on $R$ and $n_{RL}=n+1$ for one that ends on $L$. The gain is $a^{n_{RR}}b^{n_{RL}}c^{n_{LR}}d^{n_{LL}}$; substitute. On the line $\rho=(T+x)/2$ and $\lambda=(T-x)/2$ are fixed, and $a^{\rho}d^{\lambda}=a^{x}(bc)^{\lambda}\delta^{\lambda}$. On the ring $\rho$ is fixed only modulo $M$, and $(\rho+Mj,\lambda-Mj)$ multiplies the gain by $(a/d)^{Mj}$.
\end{proof}

\begin{lemma}[one site]\label{lem:fgonesite}
Let $S$ be a position product event with $m_S(x,R)=U$, $m_S(x,L)=V$, and let the functions not depend on the site. If \eqref{eq:ppcons} holds with the factor $2$, then
\begin{equation}\label{eq:fg}
F_R(T_aU+T_cV)+F_L(T_bU+T_dV)=2F_R(U)+2F_L(V).
\end{equation}
\end{lemma}
\begin{proof}
As Lemma~\ref{lem:onesite}: apply \eqref{eq:ppcons} to the walks of $S$ that end at the site $x$. A walk at $(x,R)$ with gain $g$ continues to $(x+1,R)$ with gain $ga$ and to $(x-1,L)$ with gain $gb$; one at $(x,L)$ continues with $gc$ and $gd$.
\end{proof}

By Lemma~\ref{lem:fgcoset}, on the line \eqref{eq:fg} can be asked only for $U\in\mathcal{M}(\kappa\langle\delta\rangle)$, $V\in\mathcal{M}(\kappa bd^{-1}\langle\delta\rangle)$ with one $\kappa\in H$; values on other multisets never enter $\mu$. In particular the preparation of all pairs of multisets of $A$, which Lemma~\ref{lem:prepq} provides for $D_h$, is impossible for four gains unless $\langle\delta\rangle=A$. For a subgroup $\Delta\ni\delta$ write $\mathcal{M}_R(\Delta)=\bigcup_{\kappa\in H}\mathcal{M}(\kappa\Delta)$ and $\mathcal{M}_L(\Delta)=\bigcup_{\kappa\in H}\mathcal{M}(\kappa b\Delta)$, and let $N_G$ be the set of characters $\chi$ of a group $G\ni\delta$ with $\chi(\delta)=-1$; it is empty when $r$ is odd.

\begin{theorem}[the equation for four gains]\label{thm:fgsolve}
Let $\delta\in\Delta\le A$ and $K=H\cap\Delta$. Let $F_R$ on $\mathcal{M}_R(\Delta)$ and $F_L$ on $\mathcal{M}_L(\Delta)$ be real, vanish at $\emptyset$, and satisfy \eqref{eq:fg} for all $\kappa\in H$, $U\in\mathcal{M}(\kappa\Delta)$, $V\in\mathcal{M}(\kappa bd^{-1}\Delta)$.
\begin{enumerate}[label=(\roman*)]
\item $F_R\circ T_h=F_R$ for $h\in H$, and $F_L=F_R\circ T_c$.
\item The restriction $f$ of $F_R$ to $\mathcal{M}(\Delta)$ is $K$-invariant and satisfies
\begin{equation}\label{eq:fgQ}
f(u+v)+f(u+T_\delta v)=2f(u)+2f(v)\qquad(u,v\in\mathcal{M}(\Delta)),
\end{equation}
and $(F_R,F_L)\mapsto f$ is a linear bijection onto the $K$-invariant solutions of \eqref{eq:fgQ}, with inverse $F_R(m)=f(T_{\kappa}^{-1}m)$ for $m\in\mathcal{M}(\kappa\Delta)$ and $F_L=F_R\circ T_c$.
\item These $f$ are exactly the real combinations of the sums $\sigma_\omega$, $\omega$ a character of $\Delta$ trivial on $K$, and of the products $\sigma_\chi\sigma_{\chi'}$, $\{\chi,\chi'\}\subseteq N_\Delta$ with $\chi\chi'=1$ on $K$ ($\chi=\chi'$ allowed); these functions are linearly independent, so the dimension is
\[
|\Delta/K|+\#\bigl\{\{\chi,\chi'\}\subseteq N_\Delta:\ \chi\chi'=1\text{ on }K\bigr\}.
\]
\item Conversely every such $f$ gives a solution.
\end{enumerate}
\end{theorem}
\begin{proof}
(i) The domains are closed under the translations used. With $V=\emptyset$, and with $U=\emptyset$, \eqref{eq:fg} gives $F_R(T_am)+F_L(T_bm)=2F_R(m)$ and $F_R(T_cm)+F_L(T_dm)=2F_L(m)$. Fix $m_0$ and let $\Omega$ be the finite set of pairs $(R,T_hm_0)$ and $(L,T_{hb}m_0)$, $h\in H$; put $\varphi=F$ on $\Omega$, and let $\alpha(R,m)=(R,T_am)$, $\beta(R,m)=(L,T_bm)$, $\alpha(L,m)=(R,T_cm)$, $\beta(L,m)=(L,T_dm)$. Then $\varphi=(\varphi\circ\alpha+\varphi\circ\beta)/2$, so the set where the real function $\varphi$ attains its maximum is closed under $\alpha$ and $\beta$. From $(R,m)$ one reaches $(R,T_am)$, $(R,T_{bc}m)$ and $(R,T_{bdc}m)$; these three translations have finite order and generate $H$, and from $(L,m)$ one step $\alpha$ leads to rail $R$. So $\Omega$ lies in the maximum set, $\varphi$ is constant, and $F_R(T_hm_0)=F_R(m_0)=F_L(T_bm_0)$; hence $F_L=F_R\circ T_{b^{-1}}=F_R\circ T_c$ on $\mathcal{M}_L(\Delta)$, since $bc\in H$.
(ii), (iv) Put $u=T_{\kappa}^{-1}U$ and $v=T_{\kappa a}^{-1}T_cV$; both lie in $\mathcal{M}(\Delta)$, because $\kappa bcd^{-1}=\kappa a\delta^{-1}$. Then $T_aU+T_cV=T_{\kappa a}(u+v)$ and $T_bU+T_dV=T_{\kappa b}(u+T_\delta v)$, so by (i) the equation \eqref{eq:fg} at $(U,V)$ is \eqref{eq:fgQ} at $(u,v)$, and every pair $(u,v)$ arises with $\kappa=1$. Conversely $F_R(m)=f(T_\kappa^{-1}m)$ is well defined because a non-empty $m$ in $\mathcal{M}(\kappa\Delta)\cap\mathcal{M}(\kappa'\Delta)$ has $\kappa'\kappa^{-1}\in K$; read the computation backwards.
(iii) Write $G=\Delta$; $\delta\in K$. Step 1: $E_g=\sum_{j<r}e_{g\delta^j}$ is $T_\delta$-invariant, so \eqref{eq:fgQ} gives $f(u+E_g)=f(u)+f(E_g)$. The additive function $l(m)=\frac1r\sum_g m(g)f(E_g)$ is $K$-invariant, solves \eqref{eq:fgQ} and agrees with $f$ on every $E_g$; so $G_0=f-l$ solves \eqref{eq:fgQ}, is $K$-invariant and satisfies $G_0(u+E_g)=G_0(u)$. Step 2: the all-ones multiset is a sum of the $E_g$, so $G_0$ extends to $\mathbb{Z}^G$ by $G_0(W)=G_0(W+n\mathbf 1)$, with \eqref{eq:fgQ} and invariance. Step 3: the symmetry of \eqref{eq:fgQ} gives $G_0(u+T_\delta v)=G_0(v+T_\delta u)$; replacing $v$ by $T_\delta v$ and using the invariance under $T_\delta$ ($\delta\in K$) gives $G_0(u+T_\delta^2v)=G_0(u+v)$. So $G_0$ is constant on the cosets of the subgroup $L$ generated by the $E_g$ and the $e_{g\delta^2}-e_g$. Step 4, $r$ odd: then $e_{g\delta^j}\equiv e_g$ and $re_g\equiv0$ modulo $L$, so $T_\delta$ acts trivially, \eqref{eq:fgQ} is additivity, and $rG_0(W)=G_0(rW)=0$; so $f=l$. Step 5, $r$ even: $P_g=e_g+e_{g\delta}$ is $T_\delta$-invariant modulo $L$, so $G_0(U+P_g)=G_0(U)+G_0(P_g)$, and $\frac r2P_g=E_g\equiv0$ gives $G_0(P_g)=0$. The map $\varepsilon(W)(g)=\sum_j(-1)^jW(g\delta^j)$ is onto the free group $D'=\{w:w(g\delta)=-w(g)\}$ of rank $|G|/r$, with kernel $L+\langle P_g\rangle$; so $G_0=\gamma\circ\varepsilon$, and since $T_\delta=-1$ on $D'$, $\gamma(s+t)+\gamma(s-t)=2\gamma(s)+2\gamma(t)$ on $D'$, with $\gamma$ $K$-invariant. Step 6: $\gamma(0)=0$ and $\gamma$ is even; $B(s,t)=\frac12(\gamma(s+t)-\gamma(s-t))$ satisfies $2B(s,t)+2B(s',t)=B(s+s',2t)$ and $B(s,2t)=2B(s,t)$, hence $B(s+s',2t)=2B(s+s',t)$ and $B$ is additive in $s$; it is symmetric and $B(s,s)=2\gamma(s)$, so $\gamma$ is a quadratic form, and conversely. Step 7: $f(nm)=n\,l(m)+n^2\gamma(\varepsilon m)$ determines $l$ and $\gamma\circ\varepsilon$ from $f$. The $K$-invariant additive functions on $\mathcal{M}(G)$ are spanned by the $\sigma_\omega$, $\omega$ trivial on $K$. For $\chi\in N_G$, $\sigma_\chi(\varepsilon(e_x))=r\chi(x)$, so the $|G|/r$ functionals $\sigma_\chi$ form a basis of the dual of $D'\otimes\mathbb{C}$; $T_k$ multiplies $\sigma_\chi\sigma_{\chi'}$ by $(\chi\chi')(k)$, so the invariant quadratic forms are spanned by the products with $\chi\chi'=1$ on $K$, and the real dimension equals the complex one.
\end{proof}

With $\Delta=A$ the theorem solves \eqref{eq:fg} on all multisets: the dimension is $|A/H|+\#\{\{\chi,\chi'\}\subseteq N_A:\chi\chi'=1\text{ on }H\}$, which exceeds $1+\#\{\{\chi,\bar\chi\}\subseteq N_A\}$ exactly when $H\ne A$. For example $A=\mathbb{Z}_2$, $a=d=1$, $b=c=h$ has dimension $2$, with the extra solution $F_R(m)=m(1)-m(h)$, $F_L=-F_R$; and $A=\mathbb{Z}_4$, $a=h^2$, $b=c=d=1$ has dimension $5$. On the walk these extra solutions do not survive, because the multisets that occur are confined to cosets of $\langle\delta\rangle$:

\begin{lemma}[recoding]\label{lem:fgrecode}
On the line, every walk $w$ of length $T$ from $(0,R)$ to the site $x$ has gain $\kappa_{T,x}\,g_\delta(w)$ if it ends on $R$ and $\kappa_{T,x}\,b\delta d^{-1}\,g_\delta(w)$ if it ends on $L$, where $\kappa_{T,x}=a^{x}(bc)^{(T-x)/2}$ and $g_\delta(w)$ is the gain of $w$ in $D_\delta$. Hence $m_S(x,R)=T_{\kappa_{T,x}}m^{\delta}_S(x,R)$ and $m_S(x,L)=T_{\kappa_{T,x}b\delta d^{-1}}m^{\delta}_S(x,L)$ for every event $S$ of walks of length $T$, $m^\delta$ denoting multisets in $D_\delta$.
\end{lemma}
\begin{proof}
Lemma~\ref{lem:fgcoset} for both carriers ($D_\delta$ has $a=b=c=1$ and $d=\delta$, and the same $\delta$): the ratio of the gains is $a^{\rho}(d/\delta)^{\lambda}=a^{x}(bc)^{\lambda}$ on $R$ and $a^{\rho}(d/\delta)^{\lambda-1}b$ on $L$, since $d/\delta=bc/a$.
\end{proof}

\begin{proposition}[preparation for four gains on the line]\label{prop:fgprep}
Assume the conclusion of Lemma~\ref{lem:prepq} for $D_\delta$ on the line. Then for every $\kappa\in H$, $U\in\mathcal{M}(\kappa\langle\delta\rangle)$ and $V\in\mathcal{M}(\kappa bd^{-1}\langle\delta\rangle)$ there are a site $x$ and a position product event $S$ with $m_S(x,R)=U$ and $m_S(x,L)=V$.
\end{proposition}
\begin{proof}
If $U=V=\emptyset$, an event with $A_1=\emptyset$ will do. Otherwise, since $H$ is finite and generated by $a$, $bc$ and $\delta$, there are $x,\mu_0\ge0$ with $\kappa_0=a^{x}(bc)^{\mu_0}\in\kappa\langle\delta\rangle$. Put $U'=T_{\kappa_0}^{-1}U$ and $V'=T_{\kappa_0b\delta d^{-1}}^{-1}V$, multisets of $\langle\delta\rangle$. Lemma~\ref{lem:prepq} gives an event $S'$ in $D_\delta$ of some length $T_0$, and $T_0=x+2\lambda_1$ with $\lambda_1\ge0$ because a walk of $S'$ ends at $x$; with $m^\delta_{S'}(x,R)=U'$ and $m^\delta_{S'}(x,L)=V'$. Prefixing the sets $\{-1\}$, $\{0\}$ (one step left, one step right) gives an event of length $T_0+2$ whose walks in $D_\delta$ have the same gains and the same ends, since the prefix is a single walk of gain $1$ back to $(0,R)$. Repeat it $k$ times, with $\lambda_1+k\equiv\mu_0$ modulo the order of $bc$; then $\kappa_{T_0+2k,x}=\kappa_0$, and Lemma~\ref{lem:fgrecode} gives $m_S(x,R)=U$ and $m_S(x,L)=V$.
\end{proof}

\begin{theorem}[four gains on the line]\label{thm:fgline}
Let $D$ be the line with four gains in a finite abelian group, and let $F_{(x,R)}=F_R$ and $F_{(x,L)}=F_L$ not depend on the site.
\begin{enumerate}[label=(\roman*)]
\item Given the conclusion of Lemma~\ref{lem:prepq} for $D_\delta$: if \eqref{eq:ppcons} holds with the factor $2$, then on the multisets that occur $F_R=F_L=F$ with $F(m)=\alpha|m|$ if $r$ is odd and $F(m)=\alpha|m|+\gamma|\sigma_\chi(m)|^2$, $\chi(\delta)=-1$, if $r$ is even; all such $\chi$ give the same function on the multisets that occur.
\item Conversely, on the line and on every ring, every solution of Theorem~\ref{thm:fgsolve} placed at every site satisfies $\mu(\operatorname{Ext}S)=2\mu(S)$ for every event $S$ of every length.
\end{enumerate}
\end{theorem}
\begin{proof}
(i) By Proposition~\ref{prop:fgprep} and Lemma~\ref{lem:fgonesite}, $(F_R,F_L)$ solves \eqref{eq:fg} on the domain of Theorem~\ref{thm:fgsolve} with $\Delta=\langle\delta\rangle$, so $K=\langle\delta\rangle=\Delta$. Then $|\Delta/K|=1$, the only $\omega$ is trivial and $\sigma_\omega$ is the count; the characters of $\langle\delta\rangle$ with $\chi(\delta)=-1$ exist only for $r$ even, there is exactly one, and $\chi\chi'=1$ on $K$ forces $\chi'=\chi^{-1}=\chi$; and any character of $A$ with $\chi(\delta)=-1$ restricts to it. By Theorem~\ref{thm:fgsolve}(i), (ii), $F_R(m)=f(T_\kappa^{-1}m)$ and $F_L=F_R\circ T_c$, and $|m|$ and $|\sigma_\chi(m)|^2$ are unchanged by translations. (ii) By Lemma~\ref{lem:fgcoset} the multisets $U_x=m_S(x,R)$, $V_x=m_S(x,L)$ of an event of walks of one length lie in the domain with one $\kappa$ for each $x$; \eqref{eq:fg} holds there by Theorem~\ref{thm:fgsolve}(iv) and (i), and summing over $x$ as in Theorem~\ref{thm:posprod}(iv) gives the claim; events of mixed lengths are sums of events of one length.
\end{proof}

So on the line four gains add nothing to Theorem~\ref{thm:posprodq}: the consistent rules are the count, and for $r$ even one square, with $h$ replaced by $\delta=ad(bc)^{-1}$; the linear and mixed solutions of the full equation that exist when $H\ne A$ collapse to these on the multisets that occur. On the ring $\mathbb{Z}_N$ the recoding ratio $g(w)/g_\delta(w)$ of Lemma~\ref{lem:fgrecode} changes by $(a\delta/d)^{Mj}=B^{-Mj}$ between walks to one vertex at one length whose windings differ by $j$, with $B=bc\,a^{-2}$; so it is a function of the length and the end vertex exactly when $B^{M}=1$.

\begin{proposition}[preparation for four gains on rings]\label{prop:fgprepring}
Let $D$ be the ring $\mathbb{Z}_N$ with $M\ge6$ and $B^{M}=1$, so that $\Delta_N=\langle\delta\rangle$, and assume the conclusion of Lemma~\ref{lem:prepq} for $D_\delta$ on this ring. Then the conclusion of Proposition~\ref{prop:fgprep} holds on $D$, and so does that of Theorem~\ref{thm:fgline}(i).
\end{proposition}
\begin{proof}
For a walk of length $T$ to the vertex $(\bar x,\text{rail})$ with $\rho$ right and $\lambda$ left steps put $\kappa=a^{\rho-\lambda}(bc)^{\lambda}$; the ratio is $\kappa$ on $R$ and $\kappa b\delta d^{-1}$ on $L$, as in Lemma~\ref{lem:fgrecode}. Two walks of length $T$ to $\bar x$ have $(\rho',\lambda')=(\rho+k,\lambda-k)$ with $2k\equiv0\pmod N$, so $M\mid k$, and their values of $\kappa$ differ by $B^{-k}=1$. Choose $x_0,\mu_0\ge0$ with $\kappa_0=a^{x_0}(bc)^{\mu_0}\in\kappa\langle\delta\rangle$, put $\bar x=x_0\bmod N$, and take the event $S'$ of Lemma~\ref{lem:prepq} for the translated targets, as in the proof of Proposition~\ref{prop:fgprep}; the empty target is trivial. Its common value of $\kappa$ is $a^{x_0+Ns}(bc)^{\lambda_1}$ for some integers $s$ and $\lambda_1\ge0$. Two prefixes leave the multisets of $D_\delta$ unchanged, since each is a single walk of gain $1$ from $(0,R)$ back to $(0,R)$ for $N\ge3$: one step left and one right multiplies $\kappa$ by $bc$, and $N$ steps right multiplies it by $a^{N}$. With $i\equiv-s$ modulo the order of $a^N$ prefixes of the second kind and $j\equiv\mu_0-\lambda_1$ modulo the order of $bc$ of the first, $\kappa$ becomes $\kappa_0$. The second prefix removes the parity obstruction of odd $N$, where padding by two ticks alone does not suffice.
\end{proof}

When $B^M\ne1$ the walks at one vertex differ in their winding, and the next subsection treats that case.

The four-gain question was put by the author; a builder sub-agent solved the equation and wrote Theorem~\ref{thm:fgsolve} and its proof, which the author read; the reduction of the preparation to Lemma~\ref{lem:prepq} (Lemma~\ref{lem:fgrecode} and Proposition~\ref{prop:fgprep}) is the author's. Lemma~\ref{lem:fgrecode} and the group algebra of Proposition~\ref{prop:fgprep} are checked in Lean~4 with Mathlib, and Steps~3 and~6 of Theorem~\ref{thm:fgsolve}; a challenger of the theorem found no defect in range and one fault in a side claim, since removed; a challenger of the reduction rebuilt the prepared multisets end to end in $372$ of $390$ small targets ($18$ beyond its search bound, none failing) and found two gaps and two sentences too strong, all corrected. A challenger of the ring case followed Proposition~\ref{prop:fgprepring} end to end on rings of $3$ to $8$ sites ($22{,}449$ targets passed, none failed, $19{,}971$ beyond its search bound); a builder found the obstruction on rings of $1$, $2$ and $4$ sites, which the author checked by hand. All reviewers are of the author's model family; Lemma~\ref{lem:prepq} is Director~G's and is not reproved here.

\subsection{Four arc gains on rings: the winding of the walks at one vertex}

On the ring the gain of a walk depends on two of its statistics, and Proposition~\ref{prop:fgprepring} avoided the second by assuming $B^{M}=1$. This subsection drops that assumption in two cases: one class of winding, given a property of the one-gain construction that follows from the statements of \cite{dgone}, and two classes on rings with $M\ge8$, by a preparation of three groups. It also shows that for three or more classes the preparation of all pairs does not exist. Throughout, $\lambda(w)$ is the number of left steps of a walk $w$, $\lambda_t(w)$ the number among its first $t$ steps, and
\[
\beta=B^{M}=(bc\,a^{-2})^{M},\qquad \Delta_N=\langle\delta,\beta\rangle,\qquad \mathsf{s}=\min\{j\ge1:\beta^{j}\in\langle\delta\rangle\};
\]
$\Delta_N$ is the group of the setting because $(a/d)^{M}=\beta^{-1}\delta^{-M}$. A walk is its sequence of step directions.

\begin{lemma}[gain and winding of one walk]\label{lem:fgkappa}
A walk from $(0,R)$ of length $T$ with $\lambda$ left steps has gain $\kappa\,g_\delta(w)$ if it ends on $R$ and $\kappa\,b\delta d^{-1}g_\delta(w)$ if it ends on $L$, with $\kappa=a^{T-2\lambda}(bc)^{\lambda}$. Replacing $\lambda$ by $\lambda+M$ at fixed $T$ multiplies $\kappa$ by $\beta$. On $\mathbb{Z}_N$ two walks of one length to one vertex have numbers of left steps congruent modulo $M$.
\end{lemma}
\begin{proof}
The first statement is the computation in the proof of Lemma~\ref{lem:fgrecode}, which concerns one walk and uses neither the line nor a common $\lambda$. The second is $a^{-2M}(bc)^{M}=\beta$. The site after $T$ steps is $T-2\lambda$ modulo $N$, so $2\lambda\equiv2\lambda'\pmod N$.
\end{proof}

So at a vertex $y$ of an event of one length the walks with $\lambda=l_0+jM$ have their gains in the coset $\kappa_0\beta^{j}\langle\delta\rangle$ (times $b\delta d^{-1}$ on $L$); call $j$ the \emph{level}. Levels $j$ and $j'$ give the same coset exactly when $\mathsf{s}$ divides $j-j'$.

\begin{lemma}[winding at one vertex]\label{lem:winding}
Let $S$ be a position product event of length $T$ on $\mathbb{Z}_N$, $N\ge1$, let $y$ be a vertex and $n(l)$ the number of walks of $S$ that end at $y$ with $\lambda=l$. For $l'-l\in M\mathbb{Z}$:
\begin{enumerate}[label=(\roman*)]
\item if $l'-l\ge2M$ then $n(l)\,n(l')\le n(l+M)\,n(l'-M)$;
\item $n(l)\,n(l+2M)\le n(l+M)\bigl(n(l+M)-1\bigr)$;
\item if $n(l)=1$ or $n(l+2M)=1$, then $n(l)\,n(l+2M)\le\tfrac12n(l+M)\bigl(n(l+M)-1\bigr)$.
\end{enumerate}
Hence the values of $\lambda$ at $y$ form a progression of difference $M$ without gaps, and every value that is neither the least nor the greatest is taken by at least two walks.
\end{lemma}
\begin{proof}
Let $w,w'\in S$ end at $y$ with $\lambda(w)=l$, $\lambda(w')=l'$, and $\Phi(t)=\lambda_t(w')-\lambda_t(w)$. Then $\Phi(0)=0$, $\Phi(T)\ge2M$, and $\Phi$ changes by at most one in a step; let $\tau$ be the least $t$ with $\Phi(t)=M$, so $0<\tau<T$. After $\tau$ steps the sites of $w$ and $w'$ differ by $2M\equiv0$. Let $h$ be the first $\tau$ steps of $w'$ followed by the last $T-\tau$ steps of $w$, and $h'$ the first $\tau$ steps of $w$ followed by the last $T-\tau$ steps of $w'$. The sites of $h$ are those of $w'$ up to $\tau$ and those of $w$ after it, so $h\in S$; its last step is that of $w$, so it ends at $y$; and $\lambda(h)=l+M$. Likewise $h'\in S$ ends at $y$ with $\lambda(h')=l'-M$. The map $(w,w')\mapsto(h,h')$ is injective, since $\lambda_t(h)-\lambda_t(h')=\Phi(t)$ for $t\le\tau$ recovers $\tau$ and exchanging the parts after $\tau$ recovers $(w,w')$. This is (i). If $l'=l+2M$ then $\lambda(h)=\lambda(h')$ and $h\ne h'$, because their first $\tau$ steps contain different numbers of left steps; this is (ii). For (iii) suppose two different pairs $(w_1,w_1')$, $(w_2,w_2')$ with $l'=l+2M$ have images $(h,h')$ and $(h',h)$. Their exchange times $\tau_1$, $\tau_2$ are the least $t$ with $\lambda_t(h)-\lambda_t(h')=M$ and with $\lambda_t(h')-\lambda_t(h)=M$, so they differ; let $\tau_1<\tau_2$. Exchanging back, $w_1$ begins with the first $\tau_1$ steps of $h'$ and $w_2$ with the first $\tau_2$ steps of $h$, and $\lambda_{\tau_1}(h)\ne\lambda_{\tau_1}(h')$, so $w_1\ne w_2$; in the same way $w_1'\ne w_2'$. Hence two pairs whose images are the same unordered pair have neither walk in common, and if $n(l)=1$ or $n(l+2M)=1$ the map to unordered pairs is injective. If $l<l'$ occur with $l'-l\ge2M$ then $l+M$ occurs by (i), which gives the absence of gaps by induction, and (ii) gives $n(n-1)\ge1$ at an inner value.
\end{proof}

\begin{lemma}[the two rails of one site]\label{lem:windingrails}
In the same situation let $n_R(l)$ and $n_L(l)$ be the numbers at $(x,R)$ and $(x,L)$. If $l'-l\ge M$ then $n_L(l)\,n_R(l')\le n_L(l+M)\,n_R(l'-M)$; if $l'-l\ge2M$ then $n_R(l)\,n_L(l')\le n_R(l+M)\,n_L(l'-M)$.
\end{lemma}
\begin{proof}
The same exchange. In the first case $w$ ends with a left and $w'$ with a right step, so $\Phi(T-1)=l'-l+1\ge M+1$ and $\tau<T-1$; in the second $\Phi(T-1)=l'-l-1\ge M$ and $\tau\le T-1$. In both $\tau<T$, so $h$ and $h'$ end where $w$ and $w'$ end.
\end{proof}

\begin{corollary}[no preparation of all pairs for three classes]\label{cor:nowind}
On no ring, for no position product event $S$ and no vertex $y$, does $m_S(y)$ consist of exactly three walks with gains in three different cosets of $\langle\delta\rangle$. If $\mathsf{s}\ge3$, then for $\kappa\in H$ the multiset $e_\kappa+e_{\kappa\beta}+e_{\kappa\beta^{2}}\in\mathcal{M}(\kappa\Delta_N)$ is prepared at no vertex. Whatever $\mathsf{s}$, a walk on $L$ at level $j$ together with a walk on $R$ at level $j+1$ and no other walk never occurs at one site of a position product event.
\end{corollary}
\begin{proof}
Three walks in three cosets have three levels; by Lemma~\ref{lem:winding} an inner level carries two walks. The last statement is Lemma~\ref{lem:windingrails} with $n_L(l+M)=0$.
\end{proof}

\begin{corollary}[three levels need two on the other rail]\label{cor:threetwo}
At a site of a position product event on a ring, if one rail carries walks at three or more levels and the other rail is not empty, the other rail carries walks at two or more levels. Hence, if $\mathsf{s}\ge3$, no site has on one rail a multiset that meets three cosets of $\langle\delta\rangle$ and on the other a non-empty multiset inside one coset.
\end{corollary}
\begin{proof}
The levels on the two rails of one site are congruent modulo $M$, since the site after $T$ steps is $T-2\lambda$. Let $R$ carry the levels $l_0$, $l_0+M$, $l_0+2M$ and $L$ only the level $j$. If $j\le l_0+M$, the first inequality of Lemma~\ref{lem:windingrails} with $l=j$, $l'=l_0+2M$ gives $n_L(j)\,n_R(l_0+2M)\le n_L(j+M)\,n_R(l_0+M)=0$. If $j\ge l_0+2M$, the second with $l=l_0$, $l'=j$ gives $n_R(l_0)\,n_L(j)\le n_R(l_0+M)\,n_L(j-M)=0$. With the rails exchanged: let $L$ carry $j_0$, $j_0+M$, $j_0+2M$ and $R$ only $l'$; if $l'\ge j_0+M$ the first inequality with $l=j_0$ gives $n_L(j_0)\,n_R(l')\le n_L(j_0+M)\,n_R(l'-M)=0$, and if $l'\le j_0$ the second with $l=l'$ and $j_0+2M$ in the place of $l'$ gives $n_R(l')\,n_L(j_0+2M)\le n_R(l'+M)\,n_L(j_0+M)=0$. A multiset inside one coset sits on one level when $\mathsf{s}\ge2$ (its levels are congruent modulo $\mathsf{s}M$ and without gaps), and one that meets three cosets has three levels.
\end{proof}

So for $\mathsf{s}\ge3$ the hypothesis under which Theorem~\ref{thm:fgsolve} with $\Delta=\Delta_N$ would classify the rules, the preparation of all pairs in $\mathcal{M}(\kappa\Delta_N)\times\mathcal{M}(\kappa bd^{-1}\Delta_N)$, is false; for $\mathsf{s}=2$ it is not obtained, and the last statement of the corollary excludes one pattern of levels. What can be prepared is this.

\begin{proposition}[one level]\label{prop:fgonelevel}
Let $N\ge3$ and suppose that $D_\delta$ on $\mathbb{Z}_N$ has the property $P^{+}_\delta$: for every site $x$ and all $U',V'\in\mathcal{M}(\langle\delta\rangle)$ there is a position product event with $m^\delta(x,R)=U'$, $m^\delta(x,L)=V'$ in which all walks that end at the site $x$ have one number of left steps. Then for arbitrary four gains, every $\kappa\in H$, $U\in\mathcal{M}(\kappa\langle\delta\rangle)$ and $V\in\mathcal{M}(\kappa bd^{-1}\langle\delta\rangle)$ are prepared at one site; a site-independent rule consistent with the factor $2$ is, on the multisets whose gains lie in one coset $\kappa\langle\delta\rangle$ or $\kappa b\langle\delta\rangle$, $\kappa\in H$, the function $\alpha|m|$ for $r$ odd and $\alpha|m|+\gamma|\sigma_\chi(m)|^{2}$, $\chi(\delta)=-1$, for $r$ even; and if $\mathsf{s}=1$ the conclusion of Theorem~\ref{thm:fgline}(i) holds.
\end{proposition}
\begin{proof}
As Proposition~\ref{prop:fgprepring}, with the common number of left steps in place of $B^{M}=1$: by Lemma~\ref{lem:fgkappa} all walks of the event at the site have one $\kappa=a^{x_0+Ni}(bc)^{\lambda_1}$; the prefixes ``left, right'' and ``$N$ steps right'' are single walks of $D_\delta$-gain $1$ back to $(0,R)$ for $N\ge3$, shift every $\lambda$ by $1$ and by $0$, and multiply $\kappa$ by $bc$ and $a^{N}$. Then Lemma~\ref{lem:fgonesite} and Theorem~\ref{thm:fgsolve} with $\Delta=\langle\delta\rangle$. If $\mathsf{s}=1$ every multiset that occurs lies in one coset, by Lemma~\ref{lem:fgkappa}.
\end{proof}

On rings of $1$, $2$ and $4$ sites $P^{+}_\delta$ is false for any gains: there both neighbours of a site are all the sites of the other parity, so every position product event leaves $0$ or a power of $2$ walks at a site. The construction of Lemma~\ref{lem:prepq} has the property in every replay made: $8{,}160$ requests on rings of $7$ to $24$ sites with $M\ge6$ and orders $1$ to $4$, each met with one number of left steps at the target. The reason is in the description of \cite{dgone}. There a walk with $\lambda$ left steps after $t$ steps is recorded by its \emph{deficit} $\lambda$, a mask is a set of deficits read modulo $M$, and each stage gives, for each wire, the integer deficits of its walks relative to a common origin: a table for the stage of five ticks, masks and step words for the short stages. In every stage the two candidates of an occupied entry differ from every retained entry of the next tick by at most five, by the statements of \cite{dgone}; so for $M\ge6$ a candidate is congruent to a retained entry modulo $M$ only if it equals it, and the admitted walks are those that the stage lists, with the integer deficits it lists. Hence all walks of one stock have one number of left steps, and at the last stage the stock on wire $2$ arrives on $L$ by one right and two left steps and the stock on wire $4$ on $R$ by three right steps, both at relative deficit $4$. So $P^{+}_\delta$ holds for $M\ge6$ as far as the statements of \cite{dgone} about its stages hold; that note checks the words of its five-tick stage in Lean and states the short stages by test, and they are not reproved here. The replays agree.

\begin{lemma}[invariance without preparation]\label{lem:fginv}
Let $D$ be the line or a ring with $N\ge3$, with four gains in a finite abelian group, and let $F_R$, $F_L$ not depend on the site, vanish at $\emptyset$ and satisfy \eqref{eq:ppcons} with the factor $2$. Call a multiset \emph{occurring on $R$} (on $L$) if it is $m_S(x,R)$ (resp.\ $m_S(x,L)$) for some position product event $S$ and site $x$. Then for every $m$ occurring on $R$ and every $h\in H$ the multisets $T_hm$ and $T_{hb}m$ occur on $R$ and on $L$, and $F_R(T_hm)=F_R(m)=F_L(T_{hb}m)$; in particular $F_L=F_R\circ T_c$ on the multisets occurring on $L$.
\end{lemma}
\begin{proof}
If $m=m_S(x,R)$ with $S$ of length $T\ge2$, replace $A_{T-1}$ by $A_{T-1}\cap\{x-1\}$: the walks to $(x,R)$ come from the site $x-1$ and are kept, those to $(x,L)$ come from $x+1\ne x-1$ and are removed, and then restrict $A_T$ to $\{x\}$. For $T=1$ the event is one walk. So $m$ occurs alone at its site, the other rail and all other sites being empty; likewise on $L$. For such an event \eqref{eq:ppcons} reads $F_R(T_am)+F_L(T_bm)=2F_R(m)$, and $T_am$, $T_bm$ occur alone at the sites $x+1$ and $x-1$ of the extended event with one site kept; on $L$ it reads $F_R(T_cm)+F_L(T_dm)=2F_L(m)$. This is the situation of the proof of Theorem~\ref{thm:fgsolve}(i), with the finite set $\Omega$ of the pairs $(R,T_hm_0)$, $(L,T_{hb}m_0)$, all occurring: the function is the mean of its values at two successors, the maximum set is closed under both, and every point is reached.
\end{proof}

\begin{proposition}[three groups on two levels]\label{prop:threegroups}
Let $M=N/\gcd(N,2)\ge8$ and let $\delta$ have finite order $q\ge1$. For every site $x$ and all $G_1,G_2,G_3\in\mathcal{M}(\langle\delta\rangle)$ there are position product events in $D_\delta$ on $\mathbb{Z}_N$ in which the walks that end at the site $x$ are exactly
\begin{enumerate}[label=(\roman*)]
\item on $R$ the multiset $G_2$ at one level $l$; on $L$ the multiset $G_3$ at the level $l$ and $G_1$ at the level $l+M$;
\item on $R$ the multiset $G_3$ at the level $l$ and $G_1$ at $l+M$; on $L$ the multiset $G_2$ at $l+M$;
\end{enumerate}
and no walk ends at another site. Here ``at the level $l$'' means with $\lambda=l$.
\end{proposition}

\begin{proof}
\emph{Frame.} As in \cite{dgone}, record a walk by its deficit, the number of its left steps; after $t$ steps it is at the site $t-2e$ if its deficit is $e$, so a set $E$ of deficits read modulo $M$ is the mask $\{t-2e:e\in E\}$, and a walk at deficit $e$ has the two candidates $e$ (right step) and $e+1$ (left step), of which those in $E$ modulo $M$ survive. A left step multiplies the gain by $\delta$ exactly when it follows a left step. Deficits are written relative to that of the \emph{source}, a single walk of gain $\delta^{\rho}$; a \emph{stock} is a set of walks at one relative deficit $k\in\{2,4,6\}$. At the boundaries between stages every walk has arrived by a right step. In the advances, deposits and the runs made before the source is removed the relative deficits lie in $\{0,\dots,9\}$, and the checks use that for $M\ge8$ the residues of $0,\dots,7$ are distinct, that of $8$ is $0$ or $8$ and that of $9$ is $0$, $1$ or $9$. In the runs to $M$, the merge and what follows, a stock passes the entries $7,\dots,M+2$, and the checks compare the residues of entries up to $M+3$ with those of the stocks at $2$ and $4$.

\emph{Stages.} In each line the mask is the set of the entries named, and ``refused'' lists the candidates that must not lie in it.
\begin{itemize}
\item Advance (four ticks): the source goes $0,1,2,2,2$ by $LLRR$ and every stock $k,k+1,k+1,k+2,k+2$ by $LRLR$. Refused: at tick 1 the entries $0$ and $k$ against the mask of $1$ and the $k+1$; at tick 2 the entries $1$ and $k+2\in\{4,6,8\}$ against $2$ and the $k+1\in\{3,5,7\}$; at tick 3 the entries $3$ and $k+1$ against $2$ and the $k+2$; at tick 4 the entries $3$ and $k+3\in\{5,7,9\}$ against $2$ and the $k+2\in\{4,6,8\}$. The source gains one factor $\delta$, the stocks none, all deficits grow by two, and the relative deficits are unchanged.
\item Deposit (three ticks): masks $\{0,1,2\}$, $\{0,2\}$, $\{0,2\}$, each with the stocks. The source stays at $0$ by $RRR$; a copy goes $1,2,2$ by $LLR$, with one factor $\delta$, and is then a walk of the stock at $2$. Refused: the $k+1$ at every tick, $1$ at ticks 2 and 3, $3$ at tick 3.
\item Run of a stock from $k$ to $k+1$ (two ticks, $LR$): the mask is $k+1$ with the other occupied entries, twice. Refused: $k$, then $k+2$, and $j+1$ for every other occupied entry $j$. This needs that no other packet is at $k$, $k+1$ or $k+2$ modulo $M$ and that no two others are adjacent modulo $M$. Gains are unchanged.
\item Removal of the source (one tick): the mask is the set of stocks; refused are $0$, $1$ and the $k+1$.
\item Merge of a stock at $0$ into the stock at $2$ (two ticks and an idle tick): masks $\{1\}$ and $\{2\}$ with the other stocks, then the stocks. The stock goes $1,2$ by $LL$, with one factor $\delta$; refused are $0$, then $1$, and $j+1$ for the others. On the idle tick every walk takes a right step; refused are $j+1$ for every occupied entry $j$, the merged stock included ($M+3\equiv3$).
\item Coast: the mask is the set of occupied entries; every walk takes a right step; refused are $j+1$ for every occupied entry $j$.
\item End, with stocks at $p$ and $p-2$ (two ticks): masks $\{p-1,p\}$, $\{p\}$. The stock at $p$ takes $RR$ and ends on $R$; the stock at $p-2$ takes $LL$, with one factor $\delta$, and ends on $L$ of the same site; refused are $p-2$, $p+1$, then $p-1$, $p+1$.
\end{itemize}
In each case the refused candidates are not in the mask modulo $M$: these are the comparisons of residues named above. $M\ge8$ is used in the run of the first group from $2$ to $6$, whose last refused candidate is $7$ and must not be the source at $0$, in the run of the second group, where the source at $0$ and the stock at $6$ must not be adjacent, and in the advances and deposits made with a stock at $6$, whose candidates reach $7$, $8$ and $9$. For $M=7$ the first of these fails.

\emph{Schedule.} Call the groups $A=G_1$, $B=G_2$, $C=G_3$. Deposit the walks of $A$, one deposit for each, and run the stock from $2$ to $6$; deposit $B$ and run it from $2$ to $4$; deposit $C$ at $2$. Before each deposit the source is advanced until its exponent is the one wanted, the walks of a group being taken in the cyclic order of their exponents, so at most $q-1$ advances are made in a group. Remove the source. Run $A$ from $6$ to $M\equiv0$; the entries $7,\dots,M-1,0,1$ are free of the stocks at $2$ and $4$. Merge $A$ into the stock at $2$. For mode (i) end with $p=4$: on $R$ the group $B$, on $L$ the groups $C$ and $A$. For mode (ii) run $B$ from $4$ to $M\equiv0$ and end with $p=2$: on $R$ the groups $C$ and $A$, on $L$ the group $B$. Before the end, coast for as many ticks as bring the final site to $x$. An empty group is omitted with its stages; the checks only used the absence of other packets.

\emph{Levels and gains.} The number of left steps of a walk is its deficit as an integer. With $\sigma$ the deficit of the removed source, the stock $C$ is at $\sigma+2$, $B$ at $\sigma+4$ or, after its run in mode (ii), at $\sigma+M$, and $A$ at $\sigma+M+2$ after its run and merge. The end adds $2$ to the stock that ends on $L$. So in mode (i) $B$ and $C$ end with $\lambda=\sigma+4$ and $A$ with $\sigma+4+M$; in mode (ii) $C$ ends with $\sigma+2$, and $A$ and $B$ with $\sigma+2+M$. A walk deposited when the source had exponent $\rho$ carries $\delta^{\rho+1}$ and gains one further factor at a merge and one if it ends on $L$: in mode (i) the exponents are $\rho+1$ for $B$, $\rho+2$ for $C$, $\rho+3$ for $A$; in mode (ii) $\rho+1$ for $C$, $\rho+2$ for $A$ and for $B$. The advances are chosen accordingly. No walk survives outside the two stocks named, and they end at one site.
\end{proof}

For $M\le7$ this construction fails; no other was sought. The construction and its stage checks are a builder's; the proof above is the author's reading of them, tick by tick. Replays of the whole construction by the builder ($48{,}684$ requests) and by a challenger from the description alone ($14{,}892$ requests), on rings with $M$ from $8$ to $20$, met every request \cite{comp4s}.

\begin{theorem}[two classes of winding]\label{thm:fgtwo}
Let $D$ be the ring $\mathbb{Z}_N$ with $M\ge8$ and four gains in a finite abelian group with $\mathsf{s}=2$; write $\beta^{2}=\delta^{k}$ (for $r$ even the parity of $k$ is determined). Let $F_R$, $F_L$ not depend on the site and vanish at $\emptyset$, and let \eqref{eq:ppcons} hold with the factor $2$. Then on the multisets that occur $F_L=F_R\circ T_c$, $F_R$ is invariant under translations by $H$, and on $\mathcal{M}(\Delta_N)$
\[
F_R(m)=\alpha|m|+q\bigl(\varepsilon_0^{2}+\varepsilon_1^{2}\bigr)+c\,\varepsilon_0\varepsilon_1,\qquad \varepsilon_i(m)=\sum_{j}(-1)^{j}m(\beta^{i}\delta^{j}),
\]
with $q=c=0$ if $r$ is odd and $c=0$ if $k$ is odd. These are the real combinations of the count and of the squares $|\sigma_\chi|^{2}$, $\chi$ a character of $\Delta_N$ with $\chi(\delta)=-1$, one for each pair $\{\chi,\chi^{-1}\}$: for $r$ even two squares if $k$ is even and one if $k$ is odd, for $r$ odd none. Conversely each of them is consistent on every event (Theorem~\ref{thm:fgline}(ii)).
\end{theorem}
\begin{proof}
\emph{The pairs available.} By Proposition~\ref{prop:threegroups}, Lemma~\ref{lem:fgkappa} and the two prefixes of Proposition~\ref{prop:fgonelevel}, which shift all $\lambda$ equally and leave the residues $t-2\lambda$ of a later mask unchanged, there are for every $\kappa\in H$ events with $m(x,R)=U$, $m(x,L)=V$ for all $U\in\mathcal{M}(\kappa\Delta_N)$, $V\in\mathcal{M}(\kappa bd^{-1}\Delta_N)$ such that $U$ or $V$ lies in one coset of $\langle\delta\rangle$: mode (ii) gives $U=T_{\kappa'}G_3+T_{\kappa'\beta}G_1$, $V=T_{\kappa'\beta b\delta d^{-1}}G_2$, mode (i) gives $U=T_{\kappa'}G_2$, $V=T_{\kappa'b\delta d^{-1}}(G_3+T_\beta G_1)$, and $\kappa'$ may be taken in $\kappa\langle\delta\rangle$ or in $\kappa\beta\langle\delta\rangle$, since $\beta\in H$ and $\beta^{2}\in\langle\delta\rangle$. Call these pairs \emph{admissible}. By Lemma~\ref{lem:fgonesite}, \eqref{eq:fg} holds at every admissible pair.

\emph{Reduction.} Parts (i) and (ii) of Theorem~\ref{thm:fgsolve} use \eqref{eq:fg} with $V=\emptyset$ or $U=\emptyset$, which are admissible, and a change of variables that preserves admissibility. So $F_L=F_R\circ T_c$, $F_R\circ T_h=F_R$, and the restriction $f$ of $F_R$ to $\mathcal{M}(G)$, $G=\Delta_N$, is invariant under $G$ and satisfies \eqref{eq:fgQ} whenever $u$ or $v$ lies in one coset of $\langle\delta\rangle$.

\emph{Steps 1 to 5 of Theorem~\ref{thm:fgsolve}(iii).} Step~1 uses the pair $(u,E_g)$. With $l$ as there, $G_0=f-l$ is invariant, satisfies \eqref{eq:fgQ} at admissible pairs and $G_0(u+E_g)=G_0(u)$. Step~2: $G_0(W)=G_0(W+n\mathbf{1})$ extends $G_0$ to $\mathbb{Z}^{G}$; the equation persists for $u,v\in\mathbb{Z}^{G}$ with $v$ in one coset $C$ (apply it to $u+n\mathbf{1}$ and $v+nE_C$, where $E_C$ is the sum of the elements of $C$) and with $u$ in one coset (apply it to $u+nE_C$ and $v+n\mathbf{1}$). Step~3: at $(W-e_g,e_g)$ and at $(e_g,W-e_g)$ the right sides agree, so $G_0(W-e_g+e_{g\delta})=G_0(e_g+T_\delta W-e_{g\delta})=G_0(W+e_{g\delta^{-1}}-e_g)$ by invariance; replacing $W$ by $W+e_g-e_{g\delta^{-1}}$, $G_0$ is invariant under translation by $e_{h\delta^{2}}-e_h$ for every $h$. Step~4, $r$ odd: modulo the lattice $L$ of these translations and the $E_g$, $T_\delta$ is the identity and $re_g\equiv0$; the equation for $v$ in one coset reads $G_0(u+v)=G_0(u)+G_0(v)$, and since every $W$ is a sum of two such $v$, $G_0$ is additive, and $rG_0(W)=G_0(rW)=0$. Step~5, $r$ even: $P_g=e_g+e_{g\delta}$ lies in one coset and $T_\delta P_g-P_g\in L$, so $G_0(U+P_g)=G_0(U)+G_0(P_g)$; and $E_g$ is a sum of $r/2$ translates of $P_g$ congruent to it modulo $L$, so $G_0(P_g)=0$. The $P_g$ generate the kernel of $\varepsilon=(\varepsilon_0,\varepsilon_1)\colon\mathbb{Z}^{G}\to\mathbb{Z}^{2}$, so $G_0=\gamma\circ\varepsilon$, and since $\varepsilon\circ T_\delta=-\varepsilon$,
\begin{equation}\label{eq:gammaaxis}
\gamma(s+t)+\gamma(s-t)=2\gamma(s)+2\gamma(t)\qquad\text{for all }s\in\mathbb{Z}^{2}\text{ and all }t\text{ on a coordinate axis.}
\end{equation}
\emph{In place of Step~6.} On each axis \eqref{eq:gammaaxis} is the full equation on $\mathbb{Z}$, so $\gamma(n,0)=n^{2}\gamma(1,0)$ and $\gamma(0,n)=n^{2}\gamma(0,1)$. Let $\varphi(s_0,s_1)=\gamma(s_0,s_1)-\gamma(s_0,0)-\gamma(0,s_1)$. Subtracting from \eqref{eq:gammaaxis} with $t=(t_0,0)$ the same equation at $s_1=0$ gives $\varphi(s_0+t_0,s_1)+\varphi(s_0-t_0,s_1)=2\varphi(s_0,s_1)$, so $\varphi$ is affine in $s_0$; it vanishes at $s_0=0$, so $\varphi(s_0,s_1)=s_0\varphi(1,s_1)$. In the same way $\varphi(1,s_1)=s_1\varphi(1,1)$. Hence $\gamma=q_0s_0^{2}+q_1s_1^{2}+cs_0s_1$.
\emph{Invariance under $\beta$.} $\varepsilon_1(T_\beta W)=\varepsilon_0(W)$ and $\varepsilon_0(T_\beta W)=(-1)^{k}\varepsilon_1(W)$, because $\beta^{-1}=\beta\delta^{-k}$. So $q_0=q_1$ and $c=(-1)^{k}c$. Finally $l$ is additive and invariant under $G$, so $l=\alpha|m|$. A character with $\chi(\delta)=-1$ has $\chi(\beta)^{2}=(-1)^{k}$ and $\sigma_\chi=\varepsilon_0+\chi(\beta)\varepsilon_1$; for $k$ even the two squares $(\varepsilon_0\pm\varepsilon_1)^{2}$ span the forms found, for $k$ odd $|\varepsilon_0\pm i\varepsilon_1|^{2}=\varepsilon_0^{2}+\varepsilon_1^{2}$.
\end{proof}

\begin{theorem}[three or more classes: multisets in two adjacent cosets]\label{thm:fgadj}
Let $D$ be the ring $\mathbb{Z}_N$ with $M\ge8$ and four gains in a finite abelian group with $\mathsf{s}\ge3$; write $\beta^{\mathsf{s}}=\delta^{k}$ and $C_i=\beta^{i}\langle\delta\rangle$, $i$ modulo $\mathsf{s}$. Let $F_R$, $F_L$ not depend on the site and vanish at $\emptyset$, and let \eqref{eq:ppcons} hold with the factor $2$. Then on the multisets that occur $F_L=F_R\circ T_c$ and $F_R$ is invariant under translations by $H$, and there are real $\alpha$, $q$, $c$ such that for every $i$ and every $m\in\mathcal{M}(C_i\cup C_{i+1})$
\[
F_R(m)=\alpha|m|+q\bigl(\varepsilon_i^{2}+\varepsilon_{i+1}^{2}\bigr)+c_i\,\varepsilon_i\varepsilon_{i+1},\qquad \varepsilon_j(m)=\sum_{n}(-1)^{n}m(\beta^{j}\delta^{n})\quad(0\le j<\mathsf{s}),
\]
with $c_i=c$ for $0\le i\le\mathsf{s}-2$, $c_{\mathsf{s}-1}=(-1)^{k}c$ for the pair $C_{\mathsf{s}-1},C_0$, and $q=c=0$ if $r$ is odd. All these multisets occur. For $\mathsf{s}=3$ the right side is the restriction of a combination of the count and the two squares $|\sigma_\chi|^{2}$, $\chi(\delta)=-1$, and the combination is determined by $\alpha$, $q$, $c$.
\end{theorem}
\begin{proof}
As in the proof of Theorem~\ref{thm:fgtwo}, Proposition~\ref{prop:threegroups} gives \eqref{eq:fg}, and after the reduction \eqref{eq:fgQ}, at the pairs
\[
(A_i)\colon\ u\in\mathcal{M}(C_i\cup C_{i+1}),\ v\in\mathcal{M}(C_{i+1});\qquad (B_i)\colon\ u\in\mathcal{M}(C_i),\ v\in\mathcal{M}(C_i\cup C_{i+1}),
\]
from mode (ii) and mode (i). The other two orientations, which for $\mathsf{s}=2$ come from $\beta^{2}\in\langle\delta\rangle$, are not available: by Lemma~\ref{lem:windingrails} a walk on $L$ at a level $l$ and a walk on $R$ at the level $l+M$ occur together only with walks on $L$ at $l+M$ and on $R$ at $l$; they are not used. The invariance under $H$ and $F_L=F_R\circ T_c$ on all multisets that occur are Lemma~\ref{lem:fginv}, and the change of variables of Theorem~\ref{thm:fgsolve}(ii) is as before. Fix $i=0$ and work on $\mathbb{Z}^{C_0\cup C_1}$. Step~1: for $g\in C_1$ the pair $(u,E_g)$ is of the kind $(A_0)$; for $g\in C_0$ the pair $(E_g,v)$ is of the kind $(B_0)$ and reads $2f(E_g+v)=2f(E_g)+2f(v)$, because $E_g+T_\delta v=T_\delta(E_g+v)$. Step~2 as before, with $(A_0)$ for $v$ in $C_1$ and $(B_0)$ for $u$ in $C_0$. Step~3: for $h\in C_1$ the pairs $(W-e_h,e_h)$ and $(W-e_h,e_{h\delta^{-1}})$ have equal right sides, which gives $G_0(W-e_h+e_{h\delta})=G_0(W-e_h+e_{h\delta^{-1}})$; for $h\in C_0$ the pairs $(e_h,V)$ and $(e_{h\delta^{-1}},V)$ give, with $G_0(e_h+T_\delta V)=G_0(e_{h\delta^{-1}}+V)$, that $G_0(e_h+V)=G_0(e_{h\delta^{-2}}+V)$. Steps~4 and~5 as before, the pairs $(u,P_g)$ for $g\in C_1$ and $(P_g,W)$ for $g\in C_0$. So $G_0=\gamma\circ\varepsilon$ with $\gamma$ even on $\mathbb{Z}^{2}$, and $(A_0)$ is \eqref{eq:gammaaxis} for $t$ on the second axis, while $(B_0)$, with $s$ on the first axis and $t$ arbitrary, is by evenness \eqref{eq:gammaaxis} with the two letters exchanged, that is for a step along the first axis. The argument that replaced Step~6 gives $\gamma=q_0s_0^{2}+q_0's_1^{2}+c_0s_0s_1$, and likewise on every $C_i\cup C_{i+1}$. The translation $T_\beta$ maps $\mathcal{M}(C_i\cup C_{i+1})$ onto $\mathcal{M}(C_{i+1}\cup C_{i+2})$ with $\varepsilon_{j+1}\circ T_\beta=\varepsilon_j$ for $j\le\mathsf{s}-2$ and $\varepsilon_0\circ T_\beta=(-1)^{k}\varepsilon_{\mathsf{s}-1}$; invariance gives the equalities of the coefficients. For $\mathsf{s}=3$ and $\chi(\beta)=z$, $z^{3}=(-1)^{k}$, $|\sigma_\chi|^{2}=\sum_j\varepsilon_j^{2}+2\operatorname{Re}(z)(\varepsilon_0\varepsilon_1+\varepsilon_1\varepsilon_2)+2\operatorname{Re}(z^{2})\varepsilon_0\varepsilon_2$, and the two values of $z$ up to conjugation give two independent pairs $(q,c)$.
\end{proof}

So for $\mathsf{s}=3$ a consistent rule is known, as a combination of the count and the squares, on every multiset that meets at most two cosets; what is not decided is whether it equals that combination on the multisets that meet all three.

\begin{remark}\label{rem:winding}
(a) \emph{What the theorem adds to the line.} Let $r$ be even. On the line and for $\mathsf{s}=1$ there is one square. For $\mathsf{s}=2$ and $k$ even there are two, $(\varepsilon_0+\varepsilon_1)^{2}$ and $(\varepsilon_0-\varepsilon_1)^{2}$: they add walks one winding apart with the sign $\chi(\beta)=1$ or $-1$, and a consistent rule may mix them with any real coefficients. For $k$ odd there is one square and no such freedom. The consistency requirement does not choose among the mixtures. Compare Corollary~\ref{cor:masklist}, where the quantity that position masks read is $\operatorname{Re}\chi(B)^{M}=\operatorname{Re}\chi(\beta)$. For $r$ odd only the count remains.
(b) \emph{A phase per winding that is not a rule.} For any unit $u$ the function $|\sum_w u^{\lambda(w)}\psi(w)|^{2}$ of the walks at a vertex, $\psi$ a character with $\psi(\delta)=-1$, is consistent on every event: it is the square for the carrier in which every left step carries the further gain $u$. It is a rule in the sense of this section only where it is a function of the multiset of gains. Whether it is depends on the carrier and the length. For example, on the ring of $3$ sites with $A=\mathbb{Z}_2\times\mathbb{Z}_9$, $a=(0,1)$, $b=c=0$, $d=(1,-1)$ (so $\mathsf{s}=3$), with $\psi$ the sign of the first coordinate and multisets taken up to translation by $H$ and without regard to the rail, the multiset of gains determines the function for every unit $u$ on all position product events of length at most $8$, and at length $9$ one multiset occurs with two level structures on which the functions differ unless $u^{9}=1$ \cite{comp4s}. For other gains in the same group a challenger found the first such clash at lengths $7$, $8$ or $10$.
(c) \emph{Open.} For $\mathsf{s}\ge3$: whether the count and the squares $|\sigma_\chi|^{2}$, $\chi$ a character of $\Delta_N$ with $\chi(\delta)=-1$, are the only consistent site-independent rules on the multisets that occur. They are consistent (Theorem~\ref{thm:fgline}(ii)); the proof of the converse through a preparation of all pairs is closed by Corollary~\ref{cor:nowind}, and through pairs with one side in one coset by Corollary~\ref{cor:threetwo}; Theorem~\ref{thm:fgadj} settles the multisets in two adjacent cosets. On the ring of $4$ sites with $\mathbb{Z}_2\times\mathbb{Z}_3$ and $\mathsf{s}=3$, exact ranks of the consistency equations of all position product events to length $14$ leave, on the multisets seen by length $6$, exactly the count and the two squares \cite{comp4s}; this is \textbf{COMPUTED} and supports the converse on that ring only. For $\mathsf{s}=2$: rings with $M\le7$. For $\mathsf{s}=1$: a proof of $P^{+}_\delta$ for $M\ge6$ that does not rest on \cite{dgone}, and the rings on which $P^{+}_\delta$ is false or not tested ($M\le5$).
(d) \emph{Status.} Lemmas~\ref{lem:fgkappa} to~\ref{lem:windingrails}, Lemma~\ref{lem:fginv}, Corollaries~\ref{cor:nowind} and~\ref{cor:threetwo} and Theorem~\ref{thm:fgadj} are \textbf{PROVED}; part (iii) of Lemma~\ref{lem:winding}, Corollary~\ref{cor:threetwo} and the argument of Theorem~\ref{thm:fgadj} were found by a second builder and read by the author. Proposition~\ref{prop:threegroups} and Theorem~\ref{thm:fgtwo} are \textbf{PROVED}; the proof of the proposition is a check of finitely many small comparisons, written out above and replayed. Proposition~\ref{prop:fgonelevel} is \textbf{PROVED} given $P^{+}_\delta$, and $P^{+}_\delta$ for $M\ge6$ rests on the isolation statements of \cite{dgone} and is otherwise \textbf{COMPUTED}. The step that replaces Step~6 is checked in Lean~4 with Mathlib (\texttt{WindingAxis.lean}): an equation required only along the two axes gives a quadratic form. A challenger found Theorem~\ref{thm:fgtwo} to hold by exact ranks over the rationals on $51$ carriers with separate unknowns for the two rails and multisets of up to $4$ to $7$ gains ($r$ up to $4$). The lemma on the winding was put by the author; a builder found Lemma~\ref{lem:windingrails}, disproved the preparation of four groups it had been asked for, built Proposition~\ref{prop:threegroups} on the wires of \cite{dgone} and showed that the restricted family of pairs suffices; the proof of Theorem~\ref{thm:fgtwo} as written is the author's. The reviewers are sub-agents of the author's model family; that is not independent review.
\end{remark}

\section{Gains of bounded order and a force}\label{sec:force}
